\documentclass[conference]{IEEEtran}

\usepackage{cite}
\usepackage{amsmath,amssymb,amsfonts}
\usepackage{graphicx}
\usepackage{booktabs}
\usepackage{array}
\usepackage{url}
\usepackage{xcolor}

\title{Classification Logistic Regression}

\author{
\IEEEauthorblockN{Luis Mendez}
}

\begin{document}
\maketitle

\section*{Definition}
Classification is a supervised learning task where the goal is to predict discrete class labels for given input features. Logistic regression is a classification method that models the probability of a class label as a function of the input features.

\section*{Logistic Regression}
Logistic regression predicts a categorical dependent variable based on one or more independent variables. It estimates the probability of a binary outcome using the logistic function.


\section*{Logistic Curve}
The logistic curve is an S-shaped curve that represents the relationship between the independent variables and the probability of the binary outcome. It is defined by the logistic function:
\[
\sigma(z) = \frac{1}{1 + e^{-z}}
\]
where $z$ is the linear combination of the independent variables and their coefficients.

\end{document}